\documentclass[a4paper,12pt, reqno]{amsart}

%\addtolength{\textwidth}{26mm}
%\addtolength{\textheight}{36mm}
%\addtolength{\hoffset}{-13mm}
%\addtolength{\voffset}{-18mm}

\usepackage{amsmath, amsthm, amscd, amsfonts, amssymb, graphicx, color}
\usepackage[bookmarksnumbered, colorlinks, plainpages]{hyperref}


\newcommand*{\medcap}{\mathbin{\scalebox{1.5}{\ensuremath{\cap}}}}%

\usepackage{enumitem}
\usepackage{url}
\usepackage{multirow}
\usepackage{graphicx}
\usepackage{subfigure}
\usepackage[pass]{geometry}
\usepackage{cite}
\usepackage{cancel}
\usepackage{color}
\usepackage{pgf,tikz}
\usetikzlibrary{arrows}
%\usepackage{pgf}
\usepackage{lscape}%{pdflscape}
%\usepackage{adjustbox}

% \usepackage{refcheck}
\tikzstyle{v} = [circle, draw, inner sep=2pt, minimum size=3pt, fill=black]
\usetikzlibrary{calc,shapes,decorations.pathreplacing,matrix}
\usetikzlibrary{patterns}

\tikzset{square matrix/.style={
    matrix of nodes,
    column sep=-\pgflinewidth, row sep=-\pgflinewidth,
    nodes={draw,
      minimum height=4.5pt,
      anchor=center,
      text width=4.5pt,
      align=center,
      inner sep=0pt
    },
  },
  square matrix/.default=1.2cm
}

\newtheorem{Theorem}{Theorem}[section]
\newtheorem{Definition}[Theorem]{Definition}
\newtheorem{Assumption}[Theorem]{Assumption}
\newtheorem{Lemma}[Theorem]{Lemma}
\newtheorem{Proposition}[Theorem]{Proposition}
\newtheorem{Corollary}[Theorem]{Corollary}
\newtheorem{Remark}[Theorem]{Remark}
\newtheorem{Example}[Theorem]{Example}
\newtheorem{Question}[Theorem]{Question}
\newtheorem{Problem}[Theorem]{Problem}
\newtheorem{Conjecture}[Theorem]{Conjecture}

\DeclareMathOperator{\leaf}{leaf}
\DeclareMathOperator{\diam}{diam}
\DeclareMathOperator{\cor}{cor}

\begin{document}

\title{Connected and outer-connected domination number of middle graphs}

\author[F. Kazemnejad]{Farshad Kazemnejad}
\address{Farshad Kazemnejad, Department of Mathematics, School of Sciences, Ilam University, 
P.O.Box 69315-516, Ilam, Iran.}
\email{kazemnejad.farshad@gmail.com}
\author[B. Pahlavsay]{Behnaz Pahlavsay}
\address{Behnaz Pahlavsay, Department of Mathematics, Hokkaido University, Kita 10, Nishi 8, Kita-Ku, Sapporo 060-0810, Japan.}
\email{pahlavsayb@gmail.com}
\author[E. Palezzato]{Elisa Palezzato}
\address{Elisa Palezzato, Department of Mathematics, Hokkaido University, Kita 10, Nishi 8, Kita-Ku, Sapporo 060-0810, Japan.}
\email{palezzato@math.sci.hokudai.ac.jp}
\author[M. Torielli]{Michele Torielli}
\address{Michele Torielli, Department of Mathematics, GI-CoRE GSB, Hokkaido University, Kita 10, Nishi 8, Kita-Ku, Sapporo 060-0810, Japan.}
\email{torielli@math.sci.hokudai.ac.jp}



\date{\today}


\begin{abstract}
%For a given graph $G=(V,E)$, a set $D \subseteq V(G)$ is said to be an outer-connected dominating set if $D$ is dominating and the graph $G-D$ is connected.The outer-connected domination number of a graph $G$ denoted by $\gamma{\tilde{}} _{c}(G)$ is the cardinality of a minimum outer-connected dominating set of $G$.
 In this paper, we study the notions of connected domination number and of outer-connected domination number for middle graphs. Indeed, we obtain tight bounds for this number in terms of the order of the graph $M(G)$. We also compute the outer-connected domination number of some families of graphs such as star graphs, cycle graphs, wheel graphs, complete graphs, complete bipartite graphs and some operation on graphs, explicitly. Moreover, some Nordhaus-Gaddum-like relations are presented for the outer-connected domination number of middle graphs.
\\[0.2em]

\noindent
Keywords: Connected Domination number, Outer-Connected Domination number, Domination number, Middle graph, Nordhaus-Gaddum-like relation.
\\[0.2em]

\noindent 
\end{abstract}
\maketitle
%--------------------------------------------------------------------------------
\section{\bf Introduction}

Domination problems and its many generalizations have been intensively studied in graph theory since 1950, see for example \cite{HHS5}, \cite{HHS6}, \cite{HeYe13}, \cite{farsom},  \cite{3totdominrook} and \cite{dominLatin}. In this paper, we use standard notation for graphs and we assume that every graph is non-empty, finite, undirected and simple. We refer to \cite{bondy2008graph} as a general reference on the subject.

%The concept of total domination in graphs was first introduced by Cockayne, Dawes and Hedetniemi \cite{CDH} and has been studied extensively by many researchers in the last years. The literature on this subject has been surveyed and detailed in the recent book~\cite{HeYe13}. In this paper, we study the total domination number of central graphs. In the sequel we remind some concepts and terminology which are used in this paper.

%Let $G$ be a graph with vertex set $V(G)$ of \emph{order}
%$n$ and edge set $E(G)$ of \emph{size} $m$.
%The \emph{open neighborhood} of a
%vertex $v\in V(G)$ is $N_{G}(v)=\{u\in V(G)~|~ uv\in E(G)\}$ and, similarly, the \emph{closed neighborhood} of a
%vertex $v\in V(G)$ is $N_{G}[v]=N_{G}(v)\cup \{v\}$. The \emph{degree} of a
%vertex $v\in V(G)$ is defined as $d_G(v)=\vert N_{G}(v) \vert $. 
%The \emph{distance} $d_G(v,w)$ in $G$ of two vertices $v,w\in V(G)$ is the length of the shortest path connecting $v$ and $w$.
%The \emph{diameter} of $G$, denoted $\diam(G)$, is the shortest distance between any two vertices in $G$.
%An \emph{edge cover} of a graph $G$ is a set of edges $S\subseteq E(G)$ such that every vertex of $G$ is incident to at least one edge in $S$.
%The \emph{edge cover number} of $G$, denoted by $\rho(G)$, is the minimum cardinality of an edge cover of $G$.

%\begin{Definition} 
Given a simple graph $G$, a \emph{dominating set} of $G$ is a set $S\subseteq V(G)$
such that  $N_G[v]\cap
S\neq \emptyset$, for any vertex $v\in V(G)$, where $N_G[v]$ is the closed neighborhood of $v$. The \emph{domination number} of $G$ is the minimum cardinality of a dominating set of $G$ and it is denoted by $\gamma(G)$. %Moreover, a dominating set of $G$ of cardinality $\gamma(G)$ is called a $\gamma$-set of $G$.
%\end{Definition}
%-----------------------------------

An important subclass of the dominating sets, that is central to this paper, is the class of connected dominating sets introduced in \cite{cds}.

\begin{Definition}
 A dominating set $S$ of a graph $G$ is called \emph{connected dominating set} if the induced subgraph $G[S]$ is connected. The minimum cardinality taken over all connected dominating sets in $G$ is called the \emph{connected domination number} of $G$ and is denoted by $\gamma_{c}(G)$. Moreover, a connected dominating set of $G$ of cardinality $\gamma_{c}(G)$ is called a $\gamma_{c}$-set of $G$.
\end{Definition}
%-------------------------------------------------

In \cite{ocds}, the authors, taking inspiration from the notion of connected dominating set, introduced the concept of outer-connected dominating set.
\begin{Definition}
A dominating set $S$ of a graph $G$ is called an \emph{outer-connected dominating set}  if the graph $G-S$ is connected. The minimum cardinality of an outer-connected dominating set of $G$ is called the \emph{outer-connected domination number} of $G$ and it is denoted by $\gamma{\tilde{}} _{c}(G)$.% is the cardinality of a minimum outer-connected dominating set of $G$.
\end{Definition}


%\begin{Definition} A \emph{total dominating set}, briefly TDS,  of a graph $G$ is a set $S\subseteq V(G)$
%such that  $N_G(v)\cap
%S\neq \emptyset$, for any vertex $v\in V(G)$. The \emph{total domination number} of $G$ is the minimum cardinality of a TDS of $G$ and is denoted by $\gamma_t(G)$. %Moreover, a total dominating set of $G$ of cardinality $\gamma_t(G)$ is called a $\gamma_t$-set of $G$.
%\end{Definition}
%The\emph{minimum} and \emph{maximum degree} of a vertex in $G$ are denoted by $\delta =\delta (G)$ and $\Delta =\Delta (G)$, respectively. 
%We write $K_{n}$, $C_{n}$ and $P_{n}$ for a \emph{complete} graph, a \emph{cycle} graph and a \emph{path} graph of order $n$, respectively, while $W_n$ and $K_{n_1,n_2,\cdots,n_p}$, denote a \emph{wheel} graph of order $n+1$, and a \emph{complete $p$-partite} graph, respectively. The \emph{friendship} graph $F_n$ is obtained by joining $n$ copies of the cycle graph $C_3$ with a common vertex.
%For any non-empty $S\subseteq V(G)$, we denote by $G[S]$ the subgraph of $G$ induced on the vertex set $S$. For any $v\in V(G)$, we denote by $G\setminus v$ the subgraph of $G$ induced on the vertex set $V(G)\setminus \{v\}$. Given two graphs $G$ and $H$ with distinct vertices, we can construct a new graph $G\cup H$ by imposing $V(G\cup H)=V(G)\cup V(H)$ and $E(G\cup H)=E(G)\cup E(H)$. 

%Given a graph $G$, its \emph{complement}, denoted by $\overline{G}$, is a graph with vertex set $V(G)$ such that for every two vertices $v$ and $w$, $vw\in E(\overline{G})$\ if and only if $vw\not\in E(G)$. 

%A \emph{vertex cover} of the graph $G$ is a set $D \subseteq V(G)$ such that every edge of $G$ is incident to at least one element of $D$. The \emph{vertex cover number} of $G$, denoted by $\tau(G)$, is the minimum cardinality of a vertex cover of $G$. Moreover, an \emph{edge cover} of $G$ is a set $S \subseteq E(G)$ such that every vertex of $G$ is incident to at least one edge in $S$. The \emph{edge cover number} of $G$, denoted by $\rho(G)$, is the minimum cardinality of an edge cover of $G$. An \emph{independent set} of $G$ is a subset of vertices of $G$, no two of which are adjacent. Also a \emph{maximum independent set} is an independent set of the largest cardinality in $G$. This cardinality is called the \emph{independence number} of $G$, and is denoted by $\alpha(G)$. The \emph{clique number} of $G$ is the maximum cardinality of the vertex set of a clique in $G$. For a tree graph $G$, any vertex of degree one is called a \emph{leaf} and the neighbour of a leaf is called a \emph{support vertex} of $G$. 
%The \emph{line graph} $L(G)$ of a graph $G$ is the graph with vertex set $E(G)$, where vertices $x$ and $y$ are adjacent in $L(G)$ if and only if the corresponding edges $x$ and $y$ share a common vertex in $G$.

Following our previous works \cite{DSFBME}, \cite{KPPT2021color} and \cite{KPPT2021total}, the aim of this paper is to study connected dominating sets and outer-connected dominating set of middle graphs.
The concept of middle graph of a graph was first introduced in \cite{HamYos} as an intersection graph.

\begin{Definition}
 The middle graph $M(G)$ of a graph $G$ is the graph whose vertex set is $V(G)\cup E(G)$, where two vertices $x, y$ in the vertex set of $M(G)$ are adjacent in $M(G)$ if one the following holds
 \begin{enumerate}
 \item $x, y\in E(G)$ and $x, y$ are adjacent in $G$;
 \item $x\in V (G)$, $y\in E(G)$, and $x, y$ are incident in $G$. 
 \end{enumerate}
 \end{Definition}
 
 Notice that, by definition, if $G$ is a graph of order $n$ and size $m$, then $M(G)$ is a graph of order $n+m$ and size $2m+|E(L(G))| $, where $L(G)$ is the line graph of $G$.
 
%In other words, the middle graph $M(G)$ of a graph $G$ of order $n$ and size $m$ is a graph of order $n+m$ and size $2m+|E(L(G))| $ which is obtained by subdividing each edge of $G$ exactly once and joining all the adjacent edges of $G$ in $ M(G)$. 
%It is obvious that $M(G)$ always contains the line graph $L(G)$ as an induced subgraphs.

In order to avoid confusion throughout the paper, we will use a ``standard'' notation for the vertex set and the edge set of the middle graph $M(G)$. In particular, if $V(G)=\{v_1,v_2,\dots, v_n\}$, then we set $V(M(G))=V(G)\cup \mathcal{M}$, where $\mathcal{M}=\{m_{ij}~|~ v_iv_j\in E(G)\}$ and $E(M(G))=\{v_im_{ij},v_jm_{ij}~|~ v_iv_j\in E(G)\}\cup E(L(G)) $.


%\vspace{0.2cm}
%For standard graph theory terminology not given here we refer to \cite{West}. 

The paper is organized as follows. In Section 2, we recall few known results on outer-connected domination numbers and domination numbers. In Section 3, we compute the connected domination number of the middle graph of a connected graph. In Section 4, we present some upper and lower bounds for $\gamma{\tilde{}} _{c}(M(G))$ in terms of the order of the graph $G$, we relate the outer-connected domination number of $M(G)$ to the edge cover number of $G$ and we compute explicitly $\gamma{\tilde{}} _{c}(M(G))$ for several known families of graphs. In Section 5, we compute the outer-connected domination number of the middle graphs of graphs obtained by some special operation. In Section 6, we present some Nordhaus-Gaddum like relations for the outer-connected domination number of middle graphs. We then conclude the paper with a section composed of open problems and conjectures.
%----------------------------------------------------
%\vspace{0.02cm}

\section{Preliminares}

In this short section, we recall three results which will be useful for our investigation.
%----------------------------------------------------------
%\begin{Lemma}\label{lemma:dominationisalledges}  \emph{\cite{DSFBME}}
%	Let $G$ be a graph of order $n\ge2$ without isolated vertices and $D$ a dominating set of $M(G)$. Then there exists $D'\subseteq E(G)$ a dominating set of $M(G)$ with $|D'|\le|D|$.
%\end{Lemma}
%---------------------------------------
\begin{Theorem}[\cite{ocds}]\label{outGwithdelta} 
	If $G$ is a connected graph of order $n$, then
	$$\gamma{\tilde{}} _{c}(G)\le n-\delta(G),$$
	where $\delta(G)$ is the minimum degree of a vertex in $G$.
\end{Theorem}
%--------------------------------------------------
\begin{Theorem}[\cite{DSFBME}]\label{theo:mindominconnectgraph}
	Let $G$ be a graph with $n\ge2$ vertices. Assume $G$ has no isolated vertices, then 
	$$\lceil\frac{n}{2}\rceil\le\gamma(M(G))\le n-1.$$
\end{Theorem}
%---------------------------------
\begin{Theorem}[\cite{DSFBME}]\label{covernumber}
	Let $G$ be a graph of order $n\geq 2$ with no isolated vertex. Then $$\gamma(M(G))= \rho(G),$$
	where $\rho(G)$ is the edge cover number of $G$, i.e the minimum cardinality of an edge cover of $G$.
\end{Theorem}
%----------------------------------------------------
%----------------------------------------------------------------
\section{\bf Connected domination number of middle graphs}
In this section, we calculate the exact value of the connected domination number $\gamma _{c}(M(G))$ for any connected graph $G$ of order $n \geq 3$.
%--------------------------------------------------------------------------
\begin{Theorem}\label{gammacMG}
For any connected graph $G$ of order $n \geq 3$,
$$\gamma _{c}(M(G))=n-1.$$
\end{Theorem}
%----------------------------------------------
\begin{proof}
To fix notation, let $G$ be a connected graph with vertex set $V(G)=\{v_1, \cdots, v_n\}$. Then $V((M(G)))=V(G) \cup \mathcal{M}$ where $\mathcal{M}=\{m_{ij}~|~v_iv_j \in E(G)\}$. %Let $D$ be a connected dominating set of $M(G).$ 
First assume that $G=T$ is a tree. Obviously, $\mathcal{M}$ forms a unique minimal connected path in $M(T)$ such that $N_{M(T)}[\mathcal{M}]=V(M(T))$. This implies that $\mathcal{M}$ is the minimal connected dominating set of $M(T)$ with $|\mathcal{M}|=n-1$, and hence $\gamma _{c}(M(G))=n-1$.

%\vskip 0.02cm
Now assume that $G$ is not tree. Then consider a spanning tree $H$ of $G$, and let $\mathcal{M}_1 \subseteq \mathcal{M}$ be the vertices subdividing the edges set of $H$ in $M(H)$.  Obviously, $\mathcal{M}_1$ forms a connected path in $M(H)$ such that $N_{M(H)}[\mathcal{M}_1]=V(M(H))$. Consider $\mathcal{M}_2 =\mathcal{M}\setminus\mathcal{M}_1$. Clearly $\mathcal{M}_1$  dominates all the vertices in $\mathcal{M}_2$, and hence $N_{M(G)}[\mathcal{M}_1]=\mathcal{M}_2 \cup V(M(H))=V(M(G))$. This implies that $\mathcal{M}_1$ forms the minimal connected path in $M(G)$ with $|\mathcal{M}_1|=n-1$, and hence, $\gamma _{c}(M(G))=n-1$.
\end{proof}



%----------------------------------------------------------------
\section{\bf Outer-connected domination number of middle graphs}
As we will see in this section, the computation of the outer-connected domination number is more intricate than the one for the connected domination number.

We start our study by describing a lower and an upper bound for the outer-connected domination number of the middle graph.

\begin{Theorem}\label{theo:mindomintree}
Let $G$ be a connected graph with $n\ge2$ vertices. Then $$\lceil\frac{n}{2}\rceil\le\gamma{\tilde{}} _{c}(M(G))\le n.$$
\end{Theorem}
\begin{proof} 
	%Assume that $V(G)=\{v_1, \dots, v_n\}$. Then $V(M(G))=V(G) \cup \mathcal{M}$ where $\mathcal{M}=\{m_{ij}~|~ v_iv_j\in E(G)\}$. 
	If we consider $D=V(G)$, then $D$ is an outer-connected dominating set of $M(G)$, and hence, $\gamma{\tilde{}} _{c}(M(G))\le n$,  proving the second inequality.
	
	By Theorem \ref{theo:mindominconnectgraph} we have $\gamma{\tilde{}} _{c}(M(G)) \ge \lceil\frac{n}{2}\rceil$,  proving the first inequality.
	\end{proof}
%------------------------------------------------------
As an immediate consequence of \ref{covernumber}, we have the following result.
%------------------------------------------------
\begin{Corollary}
	Let $G$ be a graph of order $n\geq 2$ with no isolated vertex. Then $$\gamma{\tilde{}} _{c}(M(G)) \ge \rho(G).$$
\end{Corollary}	
%-----------------------------------------------------------
In the next theorem, we calculate the outer-connected domination number of the middle graph of a tree.%star graph.
%------------------------------------------------------------
%\begin{Theorem}\label{theo:outminstargraph}
%	Let $ K_{1,n-1}$ be a star graph with $n\ge2$ vertices. Then $$\gamma{\tilde{}} _{c}(M(K_{1,n-1}))= n.$$
%\end{Theorem}
%%--------------------------------------
%\begin{proof}
%To fix the notation, assume $ G=K_{1,n-1}$ be a star graph with $n\ge2$ and vertex set $V(G)=\{v_0, v_1, \cdots, v_{n-1}\}$, where $d_{G}(v_0)=n-1$. Then $V(M(G))=V(G) \cup \mathcal{M}$ where $\mathcal{M}=\{m_{0i}~|~ 1 \leq i \leq n-1\}$. Let $D$ be an outer-connected dominating set of $M(G)$. Let $v_i \notin D$ for some $1 \leq i \leq n-1$. Then $m_{0i} \in D$ and so $M(G)-D$ is disconnected, which is a contradiction. Hence $\{v_1, \cdots, v_{n-1}\} \subseteq D$. On the other hand, $N_{M(G)}[v_0]\cap \{v_1, \cdots, v_{n-1}\}=\emptyset $ implies that $|D| \geq n$. By Theorem \ref{theo:mindomintree}, we have  $\gamma{\tilde{}} _{c}(M(K_{1,n-1}))= n.$
%\end{proof}
%---------------------------------------------------------------
\begin{Theorem}\label{theo:outmindomintreeneqstar}
	%Let $T\neq K_{1,n-1}$ 
	Let $T$ be a tree with $n\ge2$ vertices. Then $$\gamma{\tilde{}} _{c}(M(T))= n.$$
\end{Theorem}
%----------------------------------------------------
\begin{proof}
% To fix the notation, assume $T\neq K_{1,n-1}$ be a tree with $n\ge2$ and 
 Let $T$ be a tree of order $n$ with $V(T)=\{v_1, \dots, v_n\}$. Then $V(M(T))=V(T) \cup \mathcal{M}$ where $\mathcal{M}=\{m_{ij}~|~ v_iv_j\in E(T)\}$. %Assume that there exists $D$ an outer-connected dominating set of $M(T)$ with $|D|<n$.
 %Without loss of generality, l
 Let $\mathcal{L}=\{v_i \in V ~|~d_{T}(v_i)=1\}$ be the set of leaves of $T$ with $|\mathcal{L}|=l$, and consider
 \[
 \mathcal{M}_{1}=\{m_{ij}~|~ v_i \in \mathcal{L}~ \mbox{or}~ v_j \in \mathcal{L}\},
 \]
 
 \[
 \mathcal{M}_{2}=\mathcal{M} \setminus \mathcal{M}_{1}.
 \]
%\textbf{Claim1:} $\mathcal{L} \subseteq D$.
%\vskip 0.03cm
%\textbf{Proof of Claim1:} Because, otherwise there exists at least a vertex $v_i \in \mathcal{L}$ such that $v_i \notin D$. Then, $ N_{M(T)}[v_i]\cap D\neq \emptyset$ implies that  $ N_{M(T)}[v_i]\cap D=\{m_{ij}\}$ for some $m_{ij} \in \mathcal{M}_{1}$. So $m_{ij} \in D$ and $v_i \notin D$ implies that $M(G)-D$ is disconnected, which is a contradiction. Hence $|D \cap \mathcal{L}|=l$.

If there exists a vertex $v_i \in \mathcal{L}$ such that $v_i \notin D$, since $ N_{M(T)}[v_i]\cap D\neq \emptyset$, then $ N_{M(T)}[v_i]\cap D=\{m_{ij}\}$ for some $m_{ij} \in \mathcal{M}_{1}$. As a consequence $m_{ij} \in D$ and $v_i \notin D$, and hence $M(G)-D$ is disconnected, which is a contradiction. This implies that $\mathcal{L} \subseteq D$ and $|D \cap \mathcal{L}|=l$.

%\vskip 0.03cm
% \textbf{Claim2:} $\mathcal{M}_{2} \cap D=\emptyset$.
% \vskip 0.03cm
% \textbf{Proof of Claim2:} 
 Let $m_{ij} \in \mathcal{M}_{2}$ be such that $m_{ij} \in D$. Then, obviously  $M(G)-D$ is disconnected, which is a contradiction. As a consequence, $\mathcal{M}_{2} \cap D=\emptyset$.
 
 Now since for any $v_i \in V(T) \setminus \mathcal{L}$ we have that $ N_{M(T)}[v_i]\cap D\neq \emptyset$ and $N_{M(T)}[v_i]\cap D \subseteq \mathcal{M}_{1} \cup (V(T) \setminus \mathcal{L})$, and for every distinct $v_i, v_j \in V(T) \setminus \mathcal{L}$ we have that $(N_{M(T)}[v_i]\cap D) \cap (N_{M(T)}[v_j]\cap D)=\emptyset $, this implies that $|D \cap (\mathcal{M}_{1} \cup (V(T) \setminus \mathcal{L}))| \geq n-l$. Hence $$|D|=|D \cap \mathcal{L}|+|D \cap (\mathcal{M}_{1} \cup (V(T) \setminus \mathcal{L}))| \geq l+(n-l)=n.$$
By Theorem \ref{theo:mindomintree}, we conclude that $\gamma{\tilde{}} _{c}(M(T))= n.$


\end{proof}
%----------------------------------------------------
%By Theorems \ref{theo:outminstargraph} and   \ref{theo:outmindomintreeneqstar} , we have the following result.
%%----------------------------------------------------
%%---------------------------------------------------------------
%\begin{Corollary}\label{outcortreeandstar}
%	Let $T$ be a tree of order $n\ge2$. Then $$\gamma{\tilde{}} _{c}(M(T))= n.$$
%\end{Corollary}
%----------------------------------------------------
\begin{Remark}
By Theorem \ref{theo:outmindomintreeneqstar}, the upper bound described in Theorem \ref{theo:mindomintree} is tight.
\end{Remark}

%Note that if $T$ is a tree of order $n$, then $\gamma{\tilde{}} _{c}(M(T))=n$. Also $\gamma{\tilde{}} _{c}(T)\leq n-1$ by Theorem \ref{outGwithdelta}. So we have the following result and question. 
%------------------------------------------------------------------
\begin{Corollary}\label{gamaact<gamactildt}
If $T$ is a tree of order $n$, then
$$\gamma{\tilde{}} _{c}(T) <\gamma{\tilde{}} _{c}(M(T)). $$
\end{Corollary}
\begin{proof} By Theorem \ref{theo:outmindomintreeneqstar} $\gamma{\tilde{}} _{c}(M(T))=n$. On the other hand, $\gamma{\tilde{}} _{c}(T)\leq n-1$ by Theorem \ref{outGwithdelta}, and hence we obtain the described inequality.
\end{proof}

Recall that the \emph{line graph} $L(G)$ of a graph $G$ is the graph with vertex set $E(G)$, where vertices $x$ and $y$ are adjacent in $L(G)$ if and only if the corresponding edges $x$ and $y$ share a common vertex in $G$. Directly from this definition and Theorem \ref{theo:outmindomintreeneqstar}, we obtain the following result.

\begin{Corollary}\label{LTREE<MTREE}
For any tree $T$ of order $n\geq 2$, $$\gamma{\tilde{}} _{c}(L(T)) < \gamma{\tilde{}} _{c}(M(T)).$$
\end{Corollary}
\begin{proof}
%Assume $V=V(G)=\{v_1,\dots,v_n\}$ and  $V(M(G))=V(G) \cup \mathcal{M}$, where $\mathcal{M}=\{m_{ij}~|~ v_iv_j\in E(G)\}$. 
By definition $V(L(T))=E(T)$ and hence, $|V(L(T))|=n-1$. This clearly implies that $\gamma{\tilde{}} _{c}(L(T))\leq n-1$. Hence $\gamma{\tilde{}} _{c}(L(T)) \leq n-1 <n= \gamma{\tilde{}} _{c}(M(T))$ by Theorem \ref{theo:outmindomintreeneqstar}.
\end{proof}
%----------------------------------------------------------
By Theorem \ref{theo:outmindomintreeneqstar}, we can characterize the trees by looking at the outer-connected domination number of their middle graph.
\begin{Theorem}\label{outGisnwhenGistree}
Let $G$ be a connected graph with $n\ge4$ vertices. Then
$$\gamma{\tilde{}} _{c}(M(G))= n~~\text{if and only if $G$~ is a tree}.$$
\end{Theorem}
%-------------------------------------
\begin{proof}
Assume that $V(G)=\{v_1, \dots, v_n\}$. Then $V(M(G))=V(G) \cup \mathcal{M}$ where $\mathcal{M}=\{m_{ij}~|~ v_iv_j\in E(G)\}$. If $G$ is a tree, then $\gamma{\tilde{}} _{c}(M(G))= n$, by Theorem \ref{theo:outmindomintreeneqstar}. 
On the other hand, assume that $\gamma{\tilde{}} _{c}(M(G))= n$ and $G$ is not tree. Then $G$ contains at least a cycle of order $n \geq 3$ as a induced subgraph. Without loss of generality, assume that $G[v_1,v_2, \dots, v_k]$ is a cycle of length $k$,  for some $k \geq 3$. %Let $D$ be an outer-connected dominating set of $M(G)$. 
Consider $D=\{v_3,v_4, \dots, v_n\} \cup \{m_{12}\}$. Then $D$ is an outer-connected dominating set of $M(G)$ with $|D|=n-1$, and hence $\gamma{\tilde{}} _{c}(M(G))\leq n-1$, which is a contradiction. This implies that $G$ is a tree.
\end{proof}
%-------------------------------------------------

In the next theorem we calculate outer-connected domination number for complete graph $K_n$ where $n \ge 3$. Notice that $K_2$ is a tree and hence $\gamma{\tilde{}} _{c}(M(K_{2}))=2$ by Theorem \ref{theo:outmindomintreeneqstar}.

%------------------------------------------------
\begin{Theorem}\label{theo:outmindomincomplete}
For any complete graph $ K_{n}$ of order $n \geq 3$, we have
	  $$\gamma{\tilde{}} _{c}(M(K_{n}))=\lceil n/2\rceil$$
\end{Theorem}
%------------------------------------
\begin{proof}
Assume that $V(M(K_n))=V(K_n) \cup \mathcal{M}$ where $V(K_n)=\{v_1, \dots, v_n\}$ and $\mathcal{M}=\{m_{ij} ~|~v_i v_j \in E(G)\}$. %Let $D$ be an outer-connected dominating set of $M(G)$. 
When $n=3$, it is easy to check that $\gamma{\tilde{}} _{c}(M(K_{n}))=2$, by considering $D=\{v_1, m_{23}\}$. Now let $n \geq 4$. Assume that $n$ is even and consider \[
D=\{m_{12}, m_{34}, \dots, m_{(n-1)n}\}
.\] Then $D$ is an outer-connected dominating set of $M(G)$ with $|D|=\lceil n/2\rceil$. Similary, if $n$ is odd, consider 
\[D=\{m_{12}, m_{34}, \dots, m_{(n-2)(n-1)}, m_{(n-1)n}\}.\]

 Then $D$ is an outer-connected dominating set of $M(G)$ with $|D|=\lceil n/2\rceil$. This show that $\gamma{\tilde{}} _{c}(M(K_{n}))\leq \lceil n/2\rceil .$ On the other hand, by Theorem \ref{theo:mindomintree}, $\gamma{\tilde{}} _{c}(M(K_{n}))\geq \lceil n/2\rceil .$
 \\
%\textcolor{blue}{Changes were made by Farshad}.
\end{proof}
%------------------------------------------------
\begin{Remark}
By Theorem \ref{theo:outmindomincomplete}, the lower bound described in Theorem \ref{theo:mindomintree} is tight.
\end{Remark}
%-------------------------------------------------------


\begin{Theorem}\label{theomoutdomcycle}
For any cycle $C_n$ of order $n \geq 3$,  $$\gamma{\tilde{}} _{c}(M(C_{n}))=n-1.$$
\end{Theorem}
%------------------------------
\begin{proof}
To fix the notation, assume that $V(C_n)=\{v_1, \dots, v_n\}$ and $E(C_n)=\{v_1v_2, v_2v_3,\dots,v_{n-1}v_n,v_nv_1\}$. Then $V(M(C_n))=V(C_n)\cup \mathcal{M}$, where $\mathcal{M}=\{ m_{i(i+1)}~|~1\leq i \leq n-1 \}\cup\{m_{1n}\}$.
Consider $D=\{v_3 , \dots, v_n\} \cup \{m_{12}\}$, then $D$ is an outer-connected dominating set with $|D|=n-1$, and hence $\gamma{\tilde{}} _{c}(M(C_{n}))\le n-1.$

 Let $D$ be a minimal outer-connected dominating set of $M(C_{n})$.
%We continue our proof in the following cases:

%\vspace{0.2cm}

%\textbf{Case 1.} $|\mathcal{M} \cap D|=0$.

%\vspace{0.2cm}

If $\mathcal{M} \cap D=\emptyset$, then $N_{M(C_{n})}[v_i]\cap D \neq \emptyset$ for every $1 \leq i \leq n$, implies that $V \subseteq D$ and hence $|D| \geq n$, contradicting the minimality of $D$. 

%\vspace{0.2cm}

%\textbf{Case 2.} $|\mathcal{M} \cap D|=1$.

%\vspace{0.2cm}

Assume that $|\mathcal{M} \cap D|=1$. Without loss of generality, we can assume that $m_{12} \in D$. Then $N_{M(C_{n})}[v_i]\cap D \neq \emptyset$ for  $ i\neq 1,2$ implies that $\{v_3, \dots, v_n\} \subseteq D$ and so $\gamma{\tilde{}} _{c}(M(C_{n})) \geq n-1$, proving our statement.

%\vspace{0.2cm}

%\textbf{Case 3.} $|\mathcal{M} \cap D|=2$.

%\vspace{0.2cm}

Assume that $|\mathcal{M} \cap D|=2$. Let $m_{ij}, m_{pq} \in D$ for some $i,j,p,q$. First, assume that $m_{ij}$ is adjacent to $m_{pq}$ in $M(C_n)$. Without loss of generality, we can assume that $m_{12}, m_{23} \in D$. Since $M(C_n)-D$ is connected, then $v_2 \in D$. Moreover, $N_{M(C_{n})}[v_i]\cap D \neq \emptyset$ for  $ i\neq 1,2,3$, implies that $\{v_4, \dots, v_n\} \subseteq D$ and hence $|D| \geq (n-4+1)+3=n$, contradicting the minimality of $D$.
Assume now that $m_{ij}$ is non-adjacent to $m_{pq}$ in $M(C_n)$ with $i<j<p<q$.
Since $m_{ij}, m_{pq} \notin N_{M(C_{n})}[v_k]$ for every $k \in \{ 1, 2, \dots, n\}\setminus\{i,j,p,q\}$, then $|D \cap V(C_n)| \geq n-4$. On the other hand, since $M(C_n)-D$ is connected then $v_j, v_p \in D$. This implies that $|D|=|D \cap V(C_n)|+|D \cap \mathcal{M}| \geq (n-2)+2=n$, contradicting the minimality of $D$.

%\vspace{0.2cm}

%\textbf{Case 4.} $|\mathcal{M} \cap D|=k \geq 3$.

%\vspace{0.2cm}
Assume that $|\mathcal{M} \cap D|=k \geq 3$. Since $M(C_n)-D$ is connected, then $k<n-1$. Without loss of generality, we can assume that $\mathcal{M} \cap D=\{m_{i_1(i_1+1)},\dots,m_{i_k(i_k+1)}\}$ where $i_1<i_2<\cdots<i_k$. If $i_j+1<i_{j+1}$ for some $1\le j\le k-1$, then $M(C_n)-D$ is disconnected. This implies that $i_j+1=i_{j+1}$ for all $1\le j\le k-1$. Let $I=\{i,i+1~|~m_{i(i+1)} \in D\}$, $V_1=\{v_i \in V~|~ i \notin I\}$. $N_{M(C_{n})}[v_i]\cap D \neq \emptyset$ for  $v_i \in V_1$, implies that $V_1 \subseteq D$. Moreover, since $M(C_n)-D$ is connected, $v_{i_1+1}, \dots, v_{i_k}\in D$.
As a consequence, $D=\{m_{i_1(i_1+1)},\dots,m_{i_k(i_k+1)}\}\cup V_1\cup \{v_{i_1+1}, \dots, v_{i_k}\}$. This implies that $|D|=k+(n-k-1)+k-1=n+k-2\ge n+1$, contradicting the minimality of $D$.


%Let $I=\{i,i+1~|~m_{i(i+1)} \in D\}$. $|\mathcal{M} \cap D|=k$ implies that $k+1 \leq |I| \leq 2k$. Note $|I|=k+1$ when $M(C_n)[\mathcal{M} \cap D]$ is connected and it is isomorphic to $P_k$, and $|I|=2k$ when $V(M(C_n)[\mathcal{M} \cap D])$ is an independent set. Let $V_1=\{v_i \in V~|~ i \notin I\}$ and $V_2=\{v_i \in V~|~ i \in I\}$. $N_{M(C_{n})}[v_i]\cap D \neq \emptyset$ for  $v_i \in V_1$, implies that $V_1 \subseteq D$, and hence $|D \cap V_1| \geq n-2k$. Since $M(C_n)-D$ is connected, then there exist at least $k-1$ indices in $I$ such that $v_i \in D$ and so $|D \cap V_2| \geq k-1$. We have $$|D|=|D \cap V_1|+|D \cap V_2|+| D \cap \mathcal{M} | \geq (n-2k)+(k-1)+k=n-1$$
Therefore, $\gamma{\tilde{}} _{c}(M(C_{n}))= n-1.$
%Now since $D=\{v_3 , \cdots, v_n\} \cup \{m_{12}\}$ is an outer-connected dominating set with $|D|=n-1$, we have $\gamma{\tilde{}} _{c}(M(C_{n}))= n-1.$
\end{proof}


%------------------------------------------------------------
\begin{Theorem} \label{OUTERWN} For any wheel $W_n$ of order $n\geq 4$, 
	$$\gamma{\tilde{}} _{c}(M(W_{n}))= \lceil n/2\rceil .$$
\end{Theorem}
\begin{proof}
	To fix the notation, assume $V(W_n)=V=\{v_0,v_1,\dots, v_{n-1}\}$ and $E(W_n)=\{v_0v_1,v_0v_2,\dots, v_0v_{n-1}\}\cup\{v_1v_2, v_2v_3,\dots,v_{n-1}v_1\}$. Then we have $V(M(W_n))=V(W_n)\cup \mathcal{M}$, where $\mathcal{M}=\{ m_{0i}~|~1\leq i \leq n-1 \}\cup\{ m_{i(i+1)}~|~1\leq i \leq n-2 \}\cup\{m_{1(n-1)}\}$. %Let $D$ be an outer-connected dominating set of $M(G)$. 
	Assume that $n$ is even and consider $D=\{m_{12}, m_{34}, \dots, m_{(n-3)(n-2)}\} \cup \{m_{0(n-1)}\}$. Then $D$ is an outer-connected dominating set of $M(G)$ with $|D|=\lceil n/2\rceil$. Similarly, if $n$ is odd, consider $D=\{m_{12}, m_{34}, \dots, m_{(n-2)(n-1)}\} \cup \{m_{0(n-1)}\}$. Then $D$ is an outer-connected dominating set of $M(G)$ with $|D|=\lceil n/2\rceil$. This show that $\gamma{\tilde{}} _{c}(M(W_{n}))\leq \lceil n/2\rceil$. On the other hand, by Theorem \ref{theo:mindomintree}, $\gamma{\tilde{}} _{c}(M(W_{n}))\geq \lceil n/2\rceil .$
\end{proof}
%---------------------------------------------------
\begin{Theorem}\label{outcompletebipartitegr}
	Let $K_{n_1,n_2}$ be the complete bipartite graph with $n_2\geq n_1 \geq 2$. Then $$\gamma{\tilde{}} _{c}(M(K_{n_1,n_2}))= n_2.$$
\end{Theorem}
\begin{proof}  To fix the notation, assume $V(K_{n_1,n_2})=\{v_1,\dots, v_{n_1},u_1,\dots, u_{n_2}\}$ and $E(K_{n_1,n_2})=\{v_iu_j ~|~1\leq i \leq n_1, 1\leq j \leq n_2\}$. Then $V(M(K_{n_1,n_2}))=V(K_{n_1,n_2})\cup \mathcal{M}$, where
	$\mathcal{M}=\{ m_{ij}~|~1\leq i \leq n_1, 1\leq j \leq n_2 \}. $
Let $D$ be an outer-connected dominating set of $M(K_{n_1,n_2})$. 	
Since $D$ is a dominating set for $M(K_{n_1,n_2})$, it has to dominate $u_1,\dots, u_{n_2}$ that are all disconnected. This implies that $\gamma(M(K_{n_1,n_2}))\ge n_2$. Now since $D=\{m_{11},m_{22},\dots, m_{n_1n_1}\}\cup\{m_{n_1(n_1+1)}, m_{n_1(n_1+2)},\dots, m_{n_1n_2}\}$ is an outer-connected dominating set of $M(K_{n_1,n_2})$ with
	 $|D|=n_1+n_2-n_1=n_2$, This implies that 
	 $\gamma{\tilde{}} _{c}(M(K_{n_1,n_2}))= n_2.$
\end{proof}

%Notice that if we consider the case when $n_1=1$, the previous result gives us Proposition~\ref{prop:mindominstar}.

%\begin{Definition} The \emph{friendship} graph $F_n$ of order $2n+1$ is obtained by joining $n$ copies of the cycle graph $C_3$ with a common vertex.
%\end{Definition}

%\begin{figure}[th]
%	\centering
	%\begin{tikzpicture}
%	\draw (0,0) node[v](1){}; 
%	\draw (-1,0.5) node[v](2){}; 
%	\draw (-1,-0.5) node[v](3){};
%	\draw (1,0.5) node[v](4){};  
%	\draw (1,-0.5) node[v](5){};  
%	\draw (-0.5,1) node[v](6){};  
%	\draw (0.5,1) node[v](7){}; 
%	\draw (-0.5,-1) node[v](8){};  
%	\draw (0.5,-1) node[v](9){}; 
%	\draw (1)--(2)--(3)--(1);
%	\draw (1)--(4)--(5)--(1);
%	\draw (1)--(6)--(7)--(1);
%	\draw (1)--(8)--(9)--(1);
%	\end{tikzpicture}
	%\caption{The friendship graph $F_{4}$}\label{Fig:friendship}
%\end{figure}

\begin{Theorem}\label{outfriendship}
	Let $F_n$ be the friendship graph with $n\ge2$. Then $$\gamma{\tilde{}} _{c}(M(F_n))= n+1.$$
\end{Theorem}
\begin{proof} To fix the notation, assume $V(F_n)=\{v_0,v_1,\dots, v_{2n}\}$ and $E(F_n)=\{v_0v_1,v_0v_2,\dots, v_0v_{2n}\}\cup\{v_1v_2, v_3v_4,\dots,v_{2n-1}v_{2n}\}$. Then $V(M(F_n))=V(F_n)\cup \mathcal{M}$, where $\mathcal{M}=\{ m_i~|~1\leq i \leq 2n \}\cup\{ m_{i(i+1)}~|~1\leq i \leq 2n-1 \text{ and } i \text{ is odd}\}$. 
	

%Let $D$ be an outer-connected dominating set of $M(F_n)$. 
By Theorem~\ref{theo:mindominconnectgraph}, $\gamma(M(F_n))\ge \lceil\frac{2n+1}{2} \rceil=n+1$. Now since  $D=\{ m_{i(i+1)}~|~1\leq i \leq 2n-1 \text{ and } i \text{ is odd}\}\cup\{v_0\}$ is an outer-connected dominating set for $M(F_n)$ with $|D|=n+1$, we have $\gamma{\tilde{}} _{c}(M(F_n))= n+1.$  
\end{proof}

Putting together Theorems \ref{theomoutdomcycle} and \ref{gammacMG}, we have the following result.
%-----------------------------------------------------
\begin{Corollary}
	There exists a connected graph $G$ of order $n \geq 3$ such that
	$$\gamma _{c}(M(G))=\gamma{\tilde{}} _{c}(M(G)).$$
\end{Corollary}
%--------------------------------------------------------
\begin{Remark}
Comparing Theorem \ref{gammacMG} and Theorem \ref{theo:outmindomintreeneqstar} we conclude that for any tree $T$ we have
$$\gamma _{c}(M(T))<\gamma{\tilde{}} _{c}(M(T)).$$ 

Similarly, comparing Theorems \ref{gammacMG} and \ref{theomoutdomcycle}  we conclude that for any cycle $$\gamma _{c}(M(C_n))=\gamma{\tilde{}} _{c}(M(C_n)).$$

Finally, comparing Theorem \ref{gammacMG} and Theorems \ref{theo:outmindomincomplete} and \ref{OUTERWN} we conclude that $$\gamma{\tilde{}} _{c}(M(K_n))=\gamma{\tilde{}} _{c}(M(W_n))<\gamma _{c}(M(K_n))=\gamma _{c}(M(W_n)).$$ 

%So Theorems \ref{gammacMG}, \ref{theomoutdomcycle}, \ref{theo:outmindomincomplete}, \ref{OUTERWN} and Theorem \ref{theo:outmindomintreeneqstar} confirm the truth of the next remark.
%\end{Remark}
%--------------------------------------------------------
%\begin{Remark}
As a consequence, if $G$ be a connected graph of order $n$, then one may not conclude that
\[
\gamma _{c}(M(G)) \geq \gamma{\tilde{}} _{c}(M(G))~~~~~\mbox{or}~~~~~~\gamma _{c}(M(G)) \leq \gamma{\tilde{}} _{c}(M(G)).
\]
\end{Remark}



%--------------------------------------------
\section{Operation on graphs}
%------------------------------------
%------------------------------------------------------

In this section, we study the outer-connected domination number for the middle graph of the corona, $2$-corona and other types of graphs.
\begin{Definition}
	The \emph{corona} $G\circ K_1$ of a graph $G$ is the graph of order $2|V(G)|$ obtained from $G$ by adding a pendant edge to each vertex of $G$. The \emph{$2$-corona} $G\circ P_2$ of $G$ is the graph of order $3|V(G)|$ obtained from $G$ by attaching a path of length $2$ to each vertex of $G$ so that the resulting paths are vertex-disjoint.
\end{Definition}

\begin{Theorem}\label{theo:minoutdomincorona}
	For any connected graph $G$ of order $n\geq 2$, $$n+\lceil n/2 \rceil \leq \gamma{\tilde{}} _{c}(M(G\circ K_1)) \leq 2n.$$
\end{Theorem}
\begin{proof} 
	Assume $V(G)=\{v_1,\dots, v_{n}\}$, then $V(G\circ K_1)= \{v_{1},\dots, v_{2n}\}$ and $E(G\circ K_1)=\{v_1v_{n+1},\dots, v_nv_{2n} \}\cup E(G) $. As a consequence, $V(M(G\circ K_1))=V(G\circ K_1)\cup \mathcal{M}$, where $\mathcal{M}=\{ m_{i(n+i)}~|~1\leq i \leq n \}\cup \{ m_{ij}~|~v_iv_j\in  E(G)\}$.
	%Let $D$ be an outer-connected dominating set of $M(G\circ K_1)$. 
	Since $\{v_1,\dots, v_{2n}\}$ is an outer-connected dominating set of $M(G\circ K_1)$, we have $\gamma{\tilde{}} _{c}(M(G\circ K_1)) \leq 2n$.

	 Let $D$ be an outer-connected dominating set of $M(G\circ K_1)$. %To prove first inequality, we claim that $\{v_{n+1},\dots, v_{2n}\} \subseteq D$. 
	 Assume $v_{n+i} \notin D$ for some $1 \le i \le n$, then since $D$ is a dominating set of $M(G\circ K_1)$ this implies that $m_{i(n+i)} \in D$ and so $M(G\circ K_1)-D$ is disconnected, which is a contradiction. As a consequence $D_1=\{v_{n+1},\dots, v_{2n}\}\subseteq D$. Now since $N_{M(G\circ K_1)}[v] \cap D_{1}=\emptyset$ for all $v \in V(M(G))$, by Theorem \ref{theo:mindominconnectgraph}, we have $$\gamma{\tilde{}} _{c}(M(G\circ K_1)) \geq n+\gamma(M(G)) \geq n+\lceil n/2 \rceil.$$

\end{proof}
%---------------------------------------------------------
\begin{Remark}
The upper bound in Theorem \ref{theo:minoutdomincorona} is tight. In fact, when $G$ is a tree, then $G\circ K_1$ is also a tree and so $\gamma{\tilde{}} _{c}(M(G\circ K_1))=2n$ by Theorem \ref{theo:outmindomintreeneqstar}.
%\end{Remark}
%------------------------------------------
%\begin{Remark}

Moreover, also the lower bound in Theorem \ref{theo:minoutdomincorona} is tight. To see it, consider $G=K_n$ and 
$$D=\{v_{n+i}~|~1 \leq i \leq n\} \cup \{m_{12}, m_{34}, \dots, m_{(n-1)n}\}$$	when $n$ is even, and 
$$D=\{v_{n+i}~|~1 \leq i \leq n\} \cup \{m_{12}, m_{34}, \dots, m_{(n-2)(n-1)} m_{(n-1)n}\}$$	when $n$ is odd. In each case, $D$ is an outer-connected dominating set of $M(K_n\circ K_1)$ with $|D|=n+\lceil n/2 \rceil$.
\end{Remark}
%--------------------------------------------------------

Similarly to Theorem \ref{theo:minoutdomincorona}, we can describe lower and upper bounds for the outer-connected domination number of the middle graph of a $2$-corona graph.
\begin{Theorem}\label{theominoutdomin2corona}
	For any connected graph $G$ of order $n\geq 2$, $$2n+\lceil n/2 \rceil \leq \gamma{\tilde{}} _{c}(M(G\circ P_2)) \leq 3n.$$
\end{Theorem}
\begin{proof} 
	Assume $V(G)=\{v_1,\dots, v_n\}$, then $V(G\circ P_2)= \{v_{1},\dots, v_{3n}\}$ and $E(G\circ P_2)=\{v_iv_{n+i}, v_{n+i}v_{2n+i}~|~1\leq i \leq n \}\cup E(G) $. As a consequence, we have that $V(M(G\circ P_2))=V(G\circ P_2)\cup \mathcal{M}$, where $\mathcal{M}=\{ m_{i(n+i)},m_{(n+i)(2n+i)}~|~1\leq i \leq n \}\cup \{ m_{ij}~|~v_iv_j\in  E(G)\}$. 
	Since $\{v_1,\dots, v_{3n}\}$ is an outer-connected dominating set of $M(G\circ P_2)$, we have $\gamma{\tilde{}} _{c}(M(G\circ P_2)) \leq 3n$.
	
	 Let $D$ be an outer-connected dominating set of $M(G\circ P_2)$. To prove first inequality, we claim that $$|D_1|=|\{v_{n+i},v_{2n+i},m_{(n+i)(2n+i)}~|~1\leq i \leq n \} \cap D | \ge 2n$$
	 Assume $v_{2n+i} \notin D$ for some $1 \le i \le n$. Since $D$ is a dominating set of $M(G\circ P_2)$ this implies that $m_{(n+i)(2n+i)} \in D$ and so $M(G\circ P_2)-D$ is disconnected, which is a contradiction. Hence $\{v_{2n+i}~|~1\leq i \leq n \} \subseteq D$. Now assume  $v_{n+i} \notin D$ for some $1 \le i \le n$. Since $D$ is a dominating set of $M(G\circ P_2)$ this implies that $m_{(n+i)(2n+i)} \in D$ or $m_{i(n+i)} \in D$. If $m_{i(n+i)} \in D$, then $M(G\circ P_2)-D$ is disconnected, which is a contradiction, and hence $m_{(n+i)(2n+i)}\in D$. This shows that for every $1 \le i \le n$, we have that $v_{n+i} \in D$ or $m_{(n+i)(2n+i)}\in D$.
	 % then we conclude that $\{v_{2n+i},m_{(n+i)(2n+i)}~|~1\leq i \leq n \}=D_1 \subseteq D$ and so $|D| \ge 2n$.  
	 Now since by construction of $D_1$, we have that $N_{M(G\circ P_2)}[v] \cap D_{1}=\emptyset$ for all $v \in V(M(G))$, this implies $$\gamma{\tilde{}} _{c}(M(G\circ P_2)) \geq 2n+\gamma(M(G)) \geq 2n+\lceil n/2 \rceil$$ by Theorem \ref{theo:mindominconnectgraph}.

\end{proof}
%--------------------------------------------------
\begin{Remark}
	The upper bound in Theorem \ref{theominoutdomin2corona} is tight. This is because when $G$ is a tree, then also $G\circ P_2$ is a tree and hence $\gamma{\tilde{}} _{c}(M(G\circ P_2))=3n$ by Theorem \ref{theo:outmindomintreeneqstar}.
%\end{Remark}
%------------------------------------------
%\begin{Remark}

	Moreover, also the lower bound in Theorem \ref{theominoutdomin2corona} is tight. Consider $G=K_n$, and
	%Consider 
	$$D=\{v_{n+i},v_{2n+i}~|~1 \leq i \leq n\} \cup \{m_{12}, m_{34}, \dots, m_{(n-1)n}\}$$	when $n$ is even, and 
	$$D=\{v_{n+i},v_{2n+i}~|~1 \leq i \leq n\} \cup \{m_{12}, m_{34}, \dots, m_{(n-2)(n-1)} m_{(n-1)n}\}$$	when $n$ is odd. In both cases, $D$ is an outer-connected dominating set of $M(K_n\circ P_2)$ with $|D|=2n+\lceil n/2 \rceil$.
\end{Remark}
%-----------------------------------
%\begin{Definition}
%	The \emph{join} $G+ H$ of two graphs $G$ and $H$ is the graph with vertex set $V(G+H)=V(G)\cup V(H)$ and edge set $E(G+H)=E(G)\cup E(H)\cup \{vw~|~v\in V(G), w\in V(H)\}$.
%	%obtained by the disjoint union of $G$ and $H$ and joining each vertex of $G$ to all vertices of $H$.
%\end{Definition}

In the next two theorems, we study the outer-connected domination number of the middle graph of the join of a graph with $\overline{K_p}$. 
\begin{Theorem}\label{theomindominjoinpbig}
	For any connected graph $G$ of order $n\geq 2$ and any integer $p\geq n$, 
	 $$\gamma{\tilde{}} _{c}(M(G+\overline{K_p}))=p.$$
\end{Theorem}
\begin{proof}
	Assume $V(G)=\{v_1,\dots,v_n\}$ and $V(\overline{K_p})=\{v_{n+1},\dots,v_{n+p}\}$. Then $V(M(G+\overline{K_p}))=V(G+\overline{K_p})\cup  \mathcal{M}_1\cup  \mathcal{M}_2$ where $\mathcal{M}_1= \{m_{ij}~|~v_iv_j\in  E(G)\}$ and $\mathcal{M}_2= \{m_{i(n+j)}~|~1\leq i \leq n, 1\leq j \leq p\}$.  
	
	%Let $D$ be an outer-connected dominating set of $M(G+\overline{K_p})$. 
	By  \cite[Theorem 2.15]{DSFBME}, we have that $\gamma(M(G+\overline{K_p}))= p$, and hence $\gamma{\tilde{}} _{c}(M(G+\overline{K_p})) \geq \gamma(M(G+\overline{K_p}))= p$.
	
	On the other hand, if we consider $D=\{m_{i(n+i)}~|~1\leq i \leq n\}\cup \{m_{1(n+j)}~|~n+1\leq j \leq p\}$, then $D$ is an outer-connected dominating set of $M(G+\overline{K_p})$ with $|D|=p$, and hence $\gamma{\tilde{}} _{c}(M(G+\overline{K_p})) \le p$.
\end{proof}

\begin{Theorem}\label{theo:mindominjoinpsmallineq}
	For any connected graph $G$ of order $n\geq 2$ and any integer $p<n$, 
	$$\lceil\frac{n+p}{2}\rceil\le\gamma{\tilde{}} _{c}(M(G+\overline{K_p}))\le n.$$
\end{Theorem}
\begin{proof} The first inequality follows directly from Theorem~\ref{theo:mindomintree}.
	On the other hand, using the same notation as in the proof of Theorem~\ref{theomindominjoinpbig}, if we consider $D=\{m_{i(n+i)}~|~1\le i\le p\}\cup\{v_j~|~p+1\le j\le n\}$, then $D$ is an outer-connected dominating set of $M(G+\overline{K_p})$ with $|D|=n$, and hence we obtain the second inequality.
\end{proof}

\begin{Remark} Both inequalities in Theorem~\ref{theo:mindominjoinpsmallineq} are sharp. In fact, if we consider $G=C_4$ and $p=2$, then a direct computation shows that $\gamma{\tilde{}} _{c}(M(C_4+\overline{K_2}))=3=\lceil\frac{n+p}{2}\rceil$. Similarly, if we consider $G=C_4$ and $p=3$, then $\gamma{\tilde{}} _{c}(M(C_4+\overline{K_3}))=4=n$. 
\end{Remark}
%-------------------------------------------------
%--------------------------------------------------
\section{Nordhaus-Gaddum-like relations}

Finding a Nordhaus-Gaddum-like relation for any parameter in graph theory is one of the traditional works which started after the following theorem by Nordhaus and Gaddum from \cite{Nordhaus}.

%-----------------------------------
\begin{Theorem}[\cite{Nordhaus}]\label{Nordhaus and Gaddum}
	For any graph $G$ of order $n$, $2\sqrt{n}\leq \chi(G)+\chi(\overline{G}) \leq n+1$.
\end{Theorem}

In this section, we find Nordhaus-Gaddum-like relations for the outer-connected  domination number of middle graphs.
In particular, by Theorems \ref{theo:mindomintree} and \ref{gammacMG}, we have the following result.
%--------------------------------------------------------------
\vskip 0.02cm

%First, we calculate exact value of $\gamma _{c}(M(G))$ for any connected graph $G$ of order $n \geq 3$.
%%--------------------------------------------------------------------------
%\begin{Theorem}\label{gammacMG}
%For any connected graph $G$ of order $n \geq 3$,
%$$\gamma _{c}(M(G))=n-1.$$
%\end{Theorem}
%%----------------------------------------------
%\begin{proof}
%To fix notation, let $G$ be a connected graph of order $n \geq 3$ with vertex set $V=\{v_1, \cdots, v_n\}$. Then $V((M(G)))=V \cup \mathcal{M}$ where $\mathcal{M}=\{m_{ij}~|~v_iv_j \in E(G)\}$. Let $D$ be a connected dominating set of $M(G).$ First, let $G=T$ be a tree. Obviously, $\mathcal{M}$ forms a unique minimal connected path in $M(T)$ such that $N[\mathcal{M}]=V(M(T))$. Hence, induced subgraph by $\mathcal{M}$ is the minimal connected dominating set of $M(T)$ with $|D|=|\mathcal{M}|=\gamma _{c}(M(G))=n-1$.
%
%\vskip 0.02cm
%Now let $G$ is not tree. Then consider a spanning tree $H$ of $G$. Assume that $\mathcal{M}_1 \subseteq \mathcal{M}$ be the vertices subdividing the edges set of $H$ in $M(H)$.  Obviously, $\mathcal{M}_1$ forms a connected path in $M(H)$ such that $N[\mathcal{M}_1]=V(M(H))$. Now, adding the remaining vertices $\mathcal{M}_2 =\mathcal{M}-\mathcal{M}_1$ of $M(G)$ to $M(H)$. Since $N(m_{ij}) \cap N(m^{'}_{ij}) \neq \emptyset$ for every $m_{ij} \in \mathcal{M}_1$ and $m^{'}_{ij} \in \mathcal{M}_2$ and all vertices in $\mathcal{M}_1$ covers the vertices in $\mathcal{M}_2$, hence $N[\mathcal{M}_1]=\mathcal{M}_2 \cup V(M(H))=V(M(G))$. So induced subgraph by $\mathcal{M}_1$ forms the minimal connected path in $M(G)$ with $|D|=|\mathcal{M}_1|=\gamma _{c}(M(G))=n-1$.
%\end{proof}
%---------------------------------------------------------------
%By theorems \ref{theo:mindomintree}, \ref{gammacMG}, we have the following result.
%------------------------------------------------------
\begin{Corollary} \label{gammacandoutgamma}
	Let $G$ be a connected graph with $n\ge4$ vertices, where $G$ is not tree. Then $$n+\lceil\frac{n}{2}\rceil-1 \le \gamma _{c}(M(G))+\gamma{\tilde{}} _{c}(M(G))\le 2n-2,$$
	$$\lceil\frac{n}{2}\rceil(n-1) \le \gamma _{c}(M(G))\cdot\gamma{\tilde{}} _{c}(M(G))\le (n-1)^{2}.$$
\end{Corollary}
%---------------------------------------------------
\begin{Remark}
	The upper bounds in Corollary \ref{gammacandoutgamma} are both tight, for example when $G$ is a cycle, by Theorem \ref{theomoutdomcycle} and Theorem \ref{gammacMG}.
%\end{Remark}
%------------------------------------------------------
%\begin{Remark}

	Similarly, also the lower bounds in Corollary \ref{gammacandoutgamma} are tight, for example when $G$ is a complete graph $K_n$ or a wheel graph $W_n$, by Theorems \ref{theo:outmindomincomplete}, \ref{OUTERWN} and \ref{gammacMG}.
\end{Remark}
%----------------------------------------------------------------
%---------------------------------------------------------------------
\section{Open problems}
We conclude the paper with a series of observations and open problems related to the notion of outer-connected domination number.

%\begin{Remark}
By Corollary \ref{gamaact<gamactildt}, if $G$ is a tree of order $n$, then
$\gamma{\tilde{}} _{c}(G) <\gamma{\tilde{}} _{c}(M(G)).$ On the other hand, by Theorems \ref{theo:outmindomincomplete}, \ref{theomoutdomcycle}, \ref{OUTERWN}, \ref{outcompletebipartitegr}, \ref{outfriendship} and \cite{ocds}, it is easy to see that
\[
1=\gamma{\tilde{}} _{c}(K_n)=\gamma{\tilde{}} _{c}(W_n)<\gamma{\tilde{}} _{c}(M(K_n))=\gamma{\tilde{}} _{c}(M(W_n))=\lceil n/2\rceil,
\]
\[1=\gamma{\tilde{}} _{c}(F_n)<\gamma{\tilde{}} _{c}(M(F_n))=n+1, 
\]
\[2=\gamma{\tilde{}} _{c}(K_{n_1, n_2})<\gamma{\tilde{}} _{c}(M(K_{n_1, n_2}))=n_2 
\]
and
\[n-2=\gamma{\tilde{}} _{c}(C_n)<\gamma{\tilde{}} _{c}(M(C_n))=n-1. 
\]
%Hence, we have the following question.
These facts all support the following conjecture.
%\end{Remark}
%-------------------------------------------------------
%\begin{Question}
\begin{Conjecture}
%Whether for any connected graph $G$ of order at least $2$,$$\gamma{\tilde{}} _{c}(G) <\gamma{\tilde{}} _{c}(M(G))? $$
Let $G$ be a graph of order $n\ge 2$. Then
$$\gamma{\tilde{}} _{c}(G) <\gamma{\tilde{}} _{c}(M(G)).$$
\end{Conjecture}
%\end{Question}
%----------------------------------------
%\begin{Proposition}
%For any connected graph $G$ of order $n\geq 2$, $$\gamma{\tilde{}} _{c}(G) \leq \gamma{\tilde{}} _{c}(M(G)).$$
%\end{Proposition}
%\begin{proof}
%Assume $V=V(G)=\{v_1,\dots,v_n\}$ and  $V(M(G))=V(G) \cup \mathcal{M}$, where $\mathcal{M}=\{m_{ij}~|~ v_iv_j\in E(G)\}$.   If $G$ is a tree, then $\gamma{\tilde{}} _{c}(G) <\gamma{\tilde{}} _{c}(M(G))$, by Corollary \ref{gamaact<gamactildt}. We can then assume that $G$ is not a tree and  let $D$ be a minimal outer-connected dominating set of $M(G)$.

%\vspace{0.02cm}

%\textbf{Case1:} $\mathcal{M} \cap D=\emptyset$.
%\vskip 0.03cm
%If $\mathcal{M} \cap D=\emptyset$, then $V$, that is an independent set in $M(G)$, is a unique minimal outer-connected dominating set of $M(G)$ and so $\gamma{\tilde{}} _{c}(M(G))=n$. This implies that $G$ is a tree by Theorem \ref{outGisnwhenGistree}, which is a contradiction.

%\vspace{0.02cm}
%
%\textbf{Case2:} $\mathcal{M} \cap D\neq\emptyset$.
%\vskip 0.03cm
%We can then assume that $\mathcal{M} \cap D=\{m_{i_1j_1},\dots, m_{i_rj_r}\}\neq\emptyset$, and consider $D^{'}=(D \setminus \{m_{i_1j_1},\dots, m_{i_rj_r}\}) \cup \{v_{i_1},\dots, v_{i_r}\}$, where $\deg_{G}(v_{i_p}) \leq \deg_{G}(v_{j_p})$ for all $1\le p\le r$.  By construction, $D^{'}$ is a dominating set of $G$. %It is sufficient, we show that $G-D^{'}$ is connected. 
%Assume that $G-D^{'}$ is not connected. Without loss of generality, we can assume that in $G-D^{'}$ there is no path connecting $v_1$ and $v_2$. Clearly this implies that $v_1,v_2\notin D^{'}$ and hence that $v_1,v_2\notin D$. Since $M(G)-D$ is connected, then there exists a path in $M(G)$ connecting $v_1$ and $v_2$ such that none of the $m_{ij}$ involved are in $D$.
%****************
%\end{proof}
%--------------------------------------------------
%--------------------------------------------------
%We end this section with the following natural problem.
%-----------------------------------------
%\begin{Proposition}\label{LTREE<MTREE}
%For any tree $G$ of order $n\geq 2$, $$\gamma{\tilde{}} _{c}(L(G)) < \gamma{\tilde{}} _{c}(M(G)).$$
%\end{Proposition}
%\begin{proof}
%Assume $V=V(G)=\{v_1,\dots,v_n\}$ and  $V(M(G))=V(G) \cup \mathcal{M}$, where $\mathcal{M}=\{m_{ij}~|~ v_iv_j\in E(G)\}$. Also $V(L(G))=\mathcal{M}$ and  $|V(L(G))|=|\mathcal{M}|=n-1$. Since each subset of $n-2$ members of $\mathcal{M}$ is an outer-connected dominating set of $L(G)$, we have $\gamma{\tilde{}} _{c}(L(G))\leq n-2$. Hence $\gamma{\tilde{}} _{c}(L(G)) \leq n-2 <n= \gamma{\tilde{}} _{c}(M(G))$ by Theorem \ref{outGisnwhenGistree}.
%\end{proof}
%------------------------------------------------------------
%-------------------------------------------------------------
Similarly to the previous conjecture, it is natural to compare the outer-connected domination number of the middle graph and of the line graph.

By Corollary~\ref{LTREE<MTREE}, if $T$ is a tree, then $\gamma{\tilde{}} _{c}(L(T)) < \gamma{\tilde{}} _{c}(M(T))$. On the other hand we can obtain similar results for some known families.

%$C_n$ is isomorphic to $L(C_n)$ for every $n \geq 3$, implies that $\gamma{\tilde{}} _{c}(C_n)=\gamma{\tilde{}} _{c}(L(C_n))=n-2$ by \cite{ocds}. On the other hand $\gamma{\tilde{}} _{c}(M(C_n))=n-1$ by Theorem \ref{theomoutdomcycle}. So we have the following result.
%--------------------------------------------------------------
\begin{Proposition}\label{LCN<MCN}
For any cycle $C_n$ of order $n \geq 3$,
 $$\gamma{\tilde{}} _{c}(L(C_n)) < \gamma{\tilde{}} _{c}(M(C_n)).$$
\end{Proposition}
\begin{proof}
By definition of line graph, $C_n$ is isomorphic to $L(C_n)$ for every $n \geq 3$. This implies that $\gamma{\tilde{}} _{c}(C_n)=\gamma{\tilde{}} _{c}(L(C_n))=n-2$ by \cite{ocds}. On the other hand $\gamma{\tilde{}} _{c}(M(C_n))=n-1$ by Theorem \ref{theomoutdomcycle}, and hence $\gamma{\tilde{}} _{c}(L(C_n)) < \gamma{\tilde{}} _{c}(M(C_n)).$
\end{proof}

\begin{Proposition} \label{OUTERWNLWN} For any wheel $W_n$ of order $n\geq 5$, 
	$$\gamma{\tilde{}} _{c}(L(W_n)) < \gamma{\tilde{}} _{c}(M(W_n)).$$
\end{Proposition}
\begin{proof}
	Let $V(W_n)=\{v_0,v_1,\dots, v_{n-1}\}$ and $E(W_n)=\{v_0v_1,\dots, v_0v_{n-1}\}\cup\{v_1v_2, v_2v_3,\dots,v_{n-1}v_1\}$. Then $V(M(W_n))=V(W_n)\cup \mathcal{M}$, where $\mathcal{M}=\{ m_{0i}~|~1\leq i \leq n-1 \}\cup\{ m_{i(i+1)}~|~1\leq i \leq n-2 \}\cup\{m_{1(n-1)}\}=V(L(W_n))$ and $E(L(W_n))=\{m_{ij} m_{pq}| \{i,j\} \cap \{p,q\}=1\}$.

%\vspace{0.2cm}
	
Assume that $n$ is even and consider $D=\{m_{(2i+1)(2i+2)}| 1 \leq i \leq \lceil n/2 \rceil-2 \} \cup \{m_{01}\}$. Then $D$ is an outer-connected dominating set of $L(G)$ with $|D|=\lceil n/2\rceil-1$. Similarly, if $n$ is odd, consider $D=\{m_{0(2i+1)}|1 \leq i \leq \lceil n/2 \rceil-1 \}$. Then $D$ is an outer-connected dominating set of $L(G)$ with $|D|=\lceil n/2\rceil-1$. This show that $\gamma{\tilde{}} _{c}(M(L_{n}))\leq \lceil n/2\rceil-1$.  By Theorem \ref{OUTERWN}, $\gamma{\tilde{}} _{c}(L(W_n)) \leq \lceil n/2 \rceil-1  < \lceil n/2 \rceil= \gamma{\tilde{}} _{c}(M(W_n)).$
\end{proof}
%-----------------------------------
%----------------------------------------
\begin{Proposition}\label{LW4=MW4}
There exists a connected graph $G$ of order $n=4$ such that
$$\gamma{\tilde{}} _{c}(L(G)) = \gamma{\tilde{}} _{c}(M(G)).$$
\end{Proposition}
\begin{proof}
Consider $G=W_4$ with $V(G)=\{v_0, v_1,v_2,v_3\}$ and $E(G)=\{v_0 v_1,v_0 v_2,v_0 v_3, v_1 v_2,v_2 v_3,v_1 v_3\}.$ Then $V(M(G))=V \cup \mathcal{M}$ where $\mathcal{M}=\{m_{ij}~|~ v_iv_j\in E(G)$ and $V(L(G))= \mathcal{M}$. Assume that $D$ is a dominating set of $L(G)$ with $|D|=1$. Then there exists an index $i$ for some $1\leq i \leq 3$ such that $N_{L(G)}[m_{ij}] \cap D=\emptyset$ which is a contradiction. This implies that $\gamma(L(G)) \geq 2$, and hence that $\gamma{\tilde{}} _{c}(L(G)) \geq 2$. Now since $D=\{m_{12}, m_{03}\}$  is an outer-connected dominating set of $L(G)$ with $|D|=2$, we have $\gamma{\tilde{}} _{c}(L(G))=2$. By Theorem \ref{OUTERWN}
$\gamma{\tilde{}} _{c}(L(G)) = \gamma{\tilde{}} _{c}(M(G))=2.$
\end{proof}
%---------------------------------------
%Now, we have the following natural problem.
As a consequence of Proposition~\ref{LW4=MW4}, it is natural to ask the following
%--------------------------------------------------------
\begin{Problem}
Can we classify the graphs $G$ such that 
$$\gamma{\tilde{}} _{c}(L(G)) = \gamma{\tilde{}} _{c}(M(G))?$$
\end{Problem}
%--------------------------------------------------
%\begin{Proposition} \label{OUTERWNLWN} For any wheel $W_n$ of order $n\geq 5$, 
%	$$\gamma{\tilde{}} _{c}(L(W_n)) < \gamma{\tilde{}} _{c}(M(W_n)).$$
%\end{Proposition}
%\begin{proof}
%	To fix the notation, assume $V(W_n)=V=\{v_0,v_1,\dots, v_{n-1}\}$ and $E(W_n)=\{v_0v_1,v_0v_2,\dots, v_0v_{n-1}\}\cup\{v_1v_2, v_2v_3,\dots,v_{n-1}v_1\}$. Then we have $V(M(W_n))=V(W_n)\cup \mathcal{M}$, where $\mathcal{M}=\{ m_{0i}~|~1\leq i \leq n-1 \}\cup\{ m_{i(i+1)}~|~1\leq i \leq n-2 \}\cup\{m_{1(n-1)}\}=V(L(W_n))$ and $E(L(W_n))=\{m_{ij} m_{pq}| \{i,j\} \cap \{p,q\}=1\}$.
%
%\vspace{0.2cm}
%	
%Assume that $n$ is even and consider $D=\{m_{(2i+1)(2i+2)}| 1 \leq i \leq \lceil n/2 \rceil-2 \} \cup \{m_{01}\}$. Then $D$ is an outer-connected dominating set of $L(G)$ with $|D|=\lceil n/2\rceil-1$. Similarly, if $n$ is odd, consider $D=\{m_{0(2i+1)}|1 \leq i \leq \lceil n/2 \rceil-1 \}$. Then $D$ is an outer-connected dominating set of $L(G)$ with $|D|=\lceil n/2\rceil-1$. This show that $\gamma{\tilde{}} _{c}(M(L_{n}))\leq \lceil n/2\rceil-1$.  By Theorem \ref{OUTERWN}, $$\gamma{\tilde{}} _{c}(L(W_n)) \leq \lceil n/2 \rceil-1  < \lceil n/2 \rceil= \gamma{\tilde{}} _{c}(M(W_n)).$$
%\end{proof}
%-------------------------------------------------------------------
%\begin{Remark}\label{remarkLG<MG}
%By Proposition \ref{LTREE<MTREE}, Corollary \ref{LCN<MCN}, Proposition \ref{LW4=MW4} and Proposition \ref{OUTERWNLWN}, we have $$\gamma{\tilde{}} _{c}(L(G)) \leq \gamma{\tilde{}} _{c}(M(G)).$$
%Note in Proposition \ref{LTREE<MTREE} $|V(G)| > |E(G)|$,  in Corollary \ref{LCN<MCN} $|V(G)| = |E(G)|$ and in  Proposition \ref{OUTERWNLWN} $|V(G)| < |E(G)|$.
%\end{Remark}
%%--------------------------------------------------
%By Remark \ref{remarkLG<MG}, we end this section with the following natural problem.
%------------------------------------------------------------------
In addition, the previous results also all support the following conjecture.
\begin{Conjecture}
	%Whether for any connected graph $G$ of order at least $2$,$$\gamma{\tilde{}} _{c}(L(G)) \le \gamma{\tilde{}} _{c}(M(G))? $$
	Let $G$ be a graph of order $n\ge 2$. Then
	$$\gamma{\tilde{}} _{c}(L(G)) \le \gamma{\tilde{}} _{c}(M(G)).$$
\end{Conjecture}
%------------------------------------------------------------------------------------
\paragraph{\textbf{Acknowledgements}} During the preparation of this article the fourth author was supported by JSPS Grant-in-Aid for Early-Career Scientists (19K14493). 


%\newpage
\begin{thebibliography}{20}
\bibitem{Aouchiche13} M. Aouchiche and P. Hansen, A survey of Nordhaus-Gaddum type relations, {\em Discrete Applied Mathematics}, 161 (2013), 466--546.
\bibitem{bondy2008graph} J.~A. Bondy and U.~S.~R. Murty, {\em Graph theory},  Graduate texts in mathematics, vol. 244, Springer Science and Media, 2008.
%------------------------
\bibitem{ocds} Cyman, J. The outer-connected domination number of a graph. Australasian Journal of Combinatorics 38 (2007), 35-46.
%--------------------------------------
\bibitem{cds} E. Sampathkumar and H. B. Walikar, The connected domination number of a
graph, J. Math. Phys. Sci., 13:607-613, 1979.
%-------------------------------
%\bibitem{CM} M. Chellali and F. Maffray, Dominator colorings in some classes of graphs. Graphs Comb. 28, 97–107 (2012) 
%\bibitem{GZ} G. Chartrand and P. Zhang, Introduction to Graph Theory, McGraw-Hill, Kalamazoo, MI 2004.
%\bibitem{CDH} E. J. Cockayne, R. M. Dawes and S. T. Hedetniemi, \textit{Total domination in graphs}, Networks 10 (1980), 211-219.
\bibitem{HamYos} T. Hamada and I. Yoshimura, Traversability and connectivity of the middle graph of a graph, {\em Discrete Mathematics}, 14 (1976) 247--255.
\bibitem{HHS5}  T.~W. Haynes, S.~T. Hedetniemi and P.~J. Slater, {\em Fundamentals of Domination in Graphs}, Marcel Dekker, New York, 1998.
\bibitem{HHS6}  T.~W. Haynes, S.~T. Hedetniemi and P.~J. Slater, {\em Domination in Graphs: Advanced Topics}, Marcel Dekker, New York, 1998.
\bibitem{HeYe13} M.~A. Henning and A. Yeo, {\em Total domination in graphs}, Springer Monographs in Mathematics, 2013.
%\bibitem{G1}  R. Gera, On the dominator colorings in bipartite graphs. Inform. Technol. New Gen., ITNG07, pp. 947–952 (2007) 
%\bibitem{G2}  R. Gera, On dominator colorings in graphs. Graph Theory Notes N. Y. 52, 25–30 (2007) 
%\bibitem{GHR}  R. Gera, S. Horton and C. Rasmussen, Dominator colorings and safe clique partitions.Congr.Num.181, 19–32 (2006) 
%\bibitem{K1}  Kazemi, A.P.: Total dominator chromatic number in graphs. Manuscript, arXiv:1307.7486 (math.CO) 
%\bibitem{K2}  Kazemi, A.P.: Total dominator chromatic number and Mycieleskian graphs. Manuscript, arXiv:1307. 7706 (math.CO)
%\bibitem{Vv}  D. Michalak,  On middle and total graphs with coarseness number equal $1$, {\em Graph theory, Proc. Conf., Lag´ow/Pol. 1981, Lect. Notes Math.} 1018, Springer, Berlin 1983, 139--150.

\bibitem{DSFBME} F. Kazemnejad, B. Pahlavsay, E. Palezzato and M. Torielli, Domination number of middle graphs. To appear in Transactions on Combinatorics. https://doi.org/10.22108/TOC.2022.131151.1927

\bibitem{KPPT2021color} F. Kazemnejad, B. Pahlavsay, E. Palezzato and M. Torielli, Total dominator coloring number of middle graphs. To appear in Discrete Mathematics, Algorithms and Applications. https://doi.org/10.1142/S1793830922500768.

\bibitem{KPPT2021total} F. Kazemnejad, B. Pahlavsay, E. Palezzato and M. Torielli, Total domination number of middle graphs. {\em Electronic Journal of Graph Theory and Applications}, 10(1), 275--288, 2022. http://dx.doi.org/10.5614/ejgta.2022.10.1.19.

\bibitem{farsom} F. Kazemnejad and S. Moradi, Total Domination Number of Central Graphs, Bulletin of the Korean Mathematical Society, 56(2019), No. 4, pp.
1059-1075.

\bibitem{Nordhaus}  E.~A. Nordhaus and J.~W. Gaddum, On complementary graphs, {\em Amer. Math. Monthly}, 63 (1956), 175-177.
\bibitem{3totdominrook} B. Pahlavsay, E. Palezzato and M. Torielli, 3-tuple total domination number of rook's graphs. {\em Discussiones Mathematicae Graph Theory}. 42, 15-37, 2022. https://doi.org/10.7151/dmgt.2242.
\bibitem{dominLatin} B. Pahlavsay, E. Palezzato and M. Torielli, Domination in latin square graphs. {\em Graphs and Combinatorics}, 37(3), 971-985, 2021.
%\bibitem{West} D.~B. West, Introduction to Graph Theory, 2nd ed, prentice hall, USA, 2001.




\end{thebibliography}
 



\end{document}